%! Setting Up Determinants --- the 2 by 2 Determinant, Area, and Orientation — Unit 5, Lesson 5.0 — Exit Ticket (Answer Key)
\documentclass[10pt]{article}
\usepackage{linalg-article}
\usepackage{linalg-key}   % pulls in -boxes; do NOT also load -boxes
\newcommand{\TallMath}[1]{$\displaystyle #1\rule[-1.4em]{0pt}{3.2em}$}

\begin{document}

\pageheader{Unit 5, Lesson 5.0}{Exit Ticket --- Answer Key}
\namedateperiod

\vspace{0.15in}

No notes. Show your reasoning.

\begin{enumerate}[leftmargin=1.8em, itemsep=16pt]
  \item \textbf{Compute a determinant.} Find
        $\det\begin{bmatrix} 4 & 2 \\ 3 & 5 \end{bmatrix}$.
        \ansline{$(4)(5)-(2)(3)=20-6=14$.}
  \item \textbf{Area \& orientation.} The columns of $\begin{bmatrix} 1 & 4 \\ 2 & 3 \end{bmatrix}$ span
        a parallelogram. Find its \emph{area}, and state whether the orientation is \emph{kept} or
        \emph{flipped}. How do you know?
        \ansline{$\det=(1)(3)-(4)(2)=3-8=-5$. Area $=|-5|=5$; orientation is \emph{flipped} because the determinant is negative.}
  \item \textbf{Solve \& justify.} Find the value of $x$ that makes
        $\begin{bmatrix} 3 & x \\ 2 & 4 \end{bmatrix}$ have $\det=0$. Then explain in one sentence what
        $\det=0$ tells you about the two columns and about whether the matrix has an \emph{inverse}.
        \ansline{$(3)(4)-(x)(2)=12-2x=0\Rightarrow x=6$. When $\det=0$ the columns are parallel (the box is flat), so the matrix is \emph{not} invertible.}
\end{enumerate}

\begin{teachernote}
Sort responses into ``got it / sign slip / concept slip.'' Item~1 is a clean $ad-bc$; item~2 tests
the absolute value (area) \emph{and} reading the sign as orientation --- watch for students who give
$-5$ as an area. Item~3 is the payoff sentence: $\det=0\Rightarrow$ parallel columns $\Rightarrow$ no
inverse. If many stumble, open 5.1 with a 60-second determinant re-warm.
\end{teachernote}

\end{document}
